% PAGES: 311-313

\begin{theorem}
(i) For any square matrix $A$,
\[
q_A(x) \;=\; q_{A^t}(x),
\]
i.e., the quadratic forms associated to a square matrix and its transpose are equal.

Moreover,
\begin{equation}
q_A(x) \;=\; \frac12 x^t\left(A+A^t\right)x.
\end{equation}

(ii) If $A$ and $B$ are square matrices of the same size, then
\begin{equation}
q_A(x) = q_B(x) \quad \iff \quad A+A^t = B+B^t.
\end{equation}

(iii) For every square matrix $A$ of size $n$ there exists a unique symmetric matrix
$M$ of size $n$ such that
\[
q_A(x) = q_M(x).
\]
\end{theorem}

Note that, from (10.30), it follows that different matrices may have identical
quadratic forms.

\begin{proof}
(i) For any $x \in \R^n$, the value of the quadratic form $q_A$ evaluated at $x$ is a
number, i.e., $q_A(x) \in \R$. Then, $q_A(x) = (q_A(x))^t$, which can be written as
\[
q_A(x) \;=\; x^tAx \;=\; (q_A(x))^t \;=\; (x^tAx)^t \;=\; x^tA^tx \;=\; q_{A^t}(x).
\]

Thus, $q_A(x) = q_{A^t}(x)$ and therefore
\begin{align*}
q_A(x) \;&=\; \frac12\left(q_A(x)+q_{A^t}(x)\right) \\
&=\; \frac12\left(x^tAx+x^tA^tx\right) \\
&=\; \frac12 x^t\left(A+A^t\right)x.
\end{align*}

(ii) From (10.28), it follows that $q_A(x) = q_B(x)$ if and only if
\begin{equation}
\sum_{j=1}^n A(j,j)x_j^2 \;+\; \sum_{1\le j<k\le n} (A(j,k)+A(k,j))x_jx_k
\end{equation}
\begin{equation}
=\; \sum_{j=1}^n B(j,j)x_j^2 \;+\; \sum_{1\le j<k\le n} (B(j,k)+B(k,j))x_jx_k, \quad \forall\, x \in \R^n.
\end{equation}

By identifying the coefficients of $x_jx_k$, for all $1 \le j < k \le n$ and of
$x_j^2$, for all $j = 1:n$, from (10.31) and (10.32), we find that $q_A = q_B$ if and
only if
\begin{align}
A(j,j) &= B(j,j), \quad \forall\, 1 \le j \le n; \\
A(j,k)+A(k,j) &= B(j,k)+B(k,j), \quad \forall\, 1 \le j < k \le n.
\end{align}

Note that (10.33) and (10.34) are equivalent to saying that the matrices $A+A^t$ and
$B+B^t$ are equal. We conclude that $q_A = q_B$ if and only if $A+A^t = B+B^t$.

(iii) Let $M$ be a symmetric matrix. Then, $M^t = M$, and, using (10.30), we find
that
\begin{align*}
q_A(x) = q_M(x) \quad &\iff \quad A+A^t = M+M^t \\
&\iff \quad A+A^t = 2M \\
&\iff \quad M = \frac{A+A^t}{2}.
\end{align*}

Thus, the only symmetric matrix $M$ such that $q_A = q_M$ is $M = \dfrac{A+A^t}{2}$.
\end{proof}

\begin{theorem}
(i) The gradient and Hessian of the quadratic form $q_A : \R^n \to \R$ associated to
the $n \times n$ matrix $A$ are
\begin{align}
D(q_A(x)) &= \left((A+A^t)x\right)^t; \\
D^2(q_A(x)) &= A+A^t.
\end{align}

(ii) If $A$ is a symmetric matrix, then
\begin{align}
D(q_A(x)) &= 2(Ax)^t; \\
D^2(q_A(x)) &= 2A.
\end{align}
\end{theorem}

\begin{proof}
A proof of this result can be found in section 10.2.1; see Theorem 10.5.
\end{proof}

\section{Mathematical tools}

\subsection{Multivariable functions}

The material in this section is adapted from Stefanica \cite{36}; see section 1.6
therein for more details.

\textit{Scalar Valued Functions}

Let $f : \R^n \to \R$ be a function of $n$ variables $x_1,x_2,\ldots,x_n$ and let
$x = (x_1,x_2,\ldots,x_n)$.

If the function $f(x)$ is differentiable with respect to all the variables $x_i$,
$i = 1:n$, then the gradient of $f(x)$ is the following $1 \times n$ row vector:
\begin{equation}
Df(x) \;=\; \left(\pd{f}{x_1}(x) \quad \pd{f}{x_2}(x) \quad \ldots \quad \pd{f}{x_n}(x)\right).
\end{equation}

If the function $f(x)$ is twice differentiable with respect to all variables $x_j,
x_k$, with $1 \le j,k \le n$, then the Hessian of $f(x)$ is the following $n \times n$
matrix:
\begin{equation}
D^2f(x) \;=\; \begin{pmatrix}
\dfrac{\partial^2 f}{\partial x_1^2}(x) & \dfrac{\partial^2 f}{\partial x_2\partial x_1}(x) & \ldots & \dfrac{\partial^2 f}{\partial x_n\partial x_1}(x) \\[1.2ex]
\dfrac{\partial^2 f}{\partial x_1\partial x_2}(x) & \dfrac{\partial^2 f}{\partial x_2^2}(x) & \ldots & \dfrac{\partial^2 f}{\partial x_n\partial x_2}(x) \\[1.2ex]
\vdots & \vdots & \ddots & \vdots \\[0.6ex]
\dfrac{\partial^2 f}{\partial x_1\partial x_n}(x) & \dfrac{\partial^2 f}{\partial x_2\partial x_n}(x) & \ldots & \dfrac{\partial^2 f}{\partial x_n^2}(x)
\end{pmatrix}.
\end{equation}

In other words,
\begin{equation}
D^2f(x) \;=\; \left(\frac{\partial^2 f}{\partial x_k\partial x_j}(x)\right)_{1\le j,k\le n}.
\end{equation}

Note that, for all $j = 1:n$, it follows from (10.39) that
\[
D\left(\pd{f}{x_j}(x)\right) \;=\; \left(\frac{\partial^2 f}{\partial x_1x_j}(x) \quad \frac{\partial^2 f}{\partial x_2x_j}(x) \quad \ldots \quad \frac{\partial^2 f}{\partial x_nx_j}(x)\right),
\]
and therefore
\begin{align}
D^2f(x) \;&=\; \begin{pmatrix} D\left(\pd{f}{x_1}(x)\right) \\ D\left(\pd{f}{x_2}(x)\right) \\ \vdots \\ D\left(\pd{f}{x_n}(x)\right) \end{pmatrix}
\;=\; D\begin{pmatrix} \pd{f}{x_1}(x) \\ \pd{f}{x_2}(x) \\ \vdots \\ \pd{f}{x_n}(x) \end{pmatrix} \notag \\
&=\; D\left((Df(x))^t\right).
\end{align}

\begin{lemma}
Let $C = (c_i)_{i=1:n}$ be an $n \times 1$ column vector of real constants, and let
$x = (x_i)_{i=1:n}$ be a column vector of $n$ variables. Then,
\begin{align}
D(C^tx) &= C^t; \\
D(x^tC) &= C^t; \\
D^2(C^tx) &= \bfzero; \\
D^2(x^tC) &= \bfzero,
\end{align}
where $\bfzero$ denotes an $n \times n$ matrix with all entries equal to $0$.
\end{lemma}

\begin{proof}
Note that
\[
C^tx \;=\; x^tC \;=\; \sum_{i=1}^n c_ix_i;
\]
cf. (1.5), and let $f : \R^n \to \R$ given by
\[
f(x) \;=\; \sum_{i=1}^n c_ix_i.
\]

Then,
\[
\pd{f}{x_i}(x) = c_i, \quad \forall\, i = 1:n; \qquad \frac{\partial^2 f}{\partial x_jx_k} = 0. \quad \forall\, 1 \le j,k \le n.
\]

From (10.39) and (10.40), it follows that
\[
Df(x) \;=\; (c_1 \ c_2 \ \ldots \ c_n) \;=\; C^t; \quad D^2f(x) \;=\; \bfzero.
\qedhere
\]
\end{proof}

\begin{theorem}
(i) Let $A$ be an $n \times n$ symmetric matrix and let $x = (x_i)_{i=1:n}$ be a
column vector of $n$ variables. The gradient and Hessian of $x^tAx$ are
\begin{align}
D(x^tAx) &= 2(Ax)^t; \\
D^2(x^tAx) &= 2A.
\end{align}

% This "theorem" (Theorem 10.5) is left OPEN at the end of this file: part (ii)
% of the statement and the whole proof are on folio 298 (PDF page 314), the
% first page of the next file (sec10_c1_5.tex). That file continues directly
% with "(ii) Let A be a square matrix..." and eventually supplies
% "\end{theorem}" -- do not reopen "\begin{theorem}" or reprint "Theorem
% 10.5" there.
